\block{Conclusión}{
	Este estudio presenta una contribución importante al análisis de la economía regional de Puerto Rico al
	proponer un marco metodológico capaz de mitigar la histórica escasez de datos macroeconómicos a nivel municipal.
	A través estudio, se logró modelar de manera robusta la función de producción Cobb-Douglas
	para los 78 municipios de la isla durante el periodo 2021-2024.

	\vspace{0.5em}
	En conclusión, aunque el uso de indicadores sintéticos y probabilísticos abre una ventana crucial para entender
	las asimetrías económicas internas de Puerto Rico, se hace evidente la necesidad urgente de institucionalizar
	la recolección de datos de inversión y cuentas de producción municipales. Futuras extensiones de este trabajo
	se enfocarán en refinar el modelo de capital sintético incorporando variables espaciales explícitas (matrices
	de adyacencia municipal) para capturar efectos de desbordamiento (\textit{spillovers}) entre regiones vecinas.
}

